% ============================================================
%  课程笔记封面模板 —— 「蒙娜丽莎」风格
%  取材：达·芬奇《蒙娜丽莎》(Mona Lisa, c.1503–1519)：晕涂法暖褐调、
%        蓝灰远景山水、交叠的双手、含蓄的微笑。
%  与 cover_Impression风.tex 同构（眉题 / 主标题 / 菱形饰线 /
%  副标题 / 底部作者日期），直接 \input{} 复用：
%      % ============================================================
%  课程笔记封面模板 —— 「蒙娜丽莎」风格
%  取材：达·芬奇《蒙娜丽莎》(Mona Lisa, c.1503–1519)：晕涂法暖褐调、
%        蓝灰远景山水、交叠的双手、含蓄的微笑。
%  与 cover_Impression风.tex 同构（眉题 / 主标题 / 菱形饰线 /
%  副标题 / 底部作者日期），直接 \input{} 复用：
%      % ============================================================
%  课程笔记封面模板 —— 「蒙娜丽莎」风格
%  取材：达·芬奇《蒙娜丽莎》(Mona Lisa, c.1503–1519)：晕涂法暖褐调、
%        蓝灰远景山水、交叠的双手、含蓄的微笑。
%  与 cover_Impression风.tex 同构（眉题 / 主标题 / 菱形饰线 /
%  副标题 / 底部作者日期），直接 \input{} 复用：
%      % ============================================================
%  课程笔记封面模板 —— 「蒙娜丽莎」风格
%  取材：达·芬奇《蒙娜丽莎》(Mona Lisa, c.1503–1519)：晕涂法暖褐调、
%        蓝灰远景山水、交叠的双手、含蓄的微笑。
%  与 cover_Impression风.tex 同构（眉题 / 主标题 / 菱形饰线 /
%  副标题 / 底部作者日期），直接 \input{} 复用：
%      \input{cover_蒙娜丽莎风}
%
%  配色：ml* 前缀（达·芬奇取色）。本文件用 \providecolor 兜底，
%        单独复制到任何笔记本也能编译；想整体换调改这里即可。
%
%  注意：整幅画面绘于 \begin{scope}[shift={(current page.south west)}] 内，
%        所有坐标以「页面左下角为原点、单位 cm」书写；脱离该 scope
%        直接写 (x,y)，图形会落到页面外而不可见。
%
%  可用字段（在主文件导言区用 \newcommand 预定义即可覆盖缺省值）：
%      \notename     主标题（必填，主模板已定义）
%      \noteauthor   作者（必填，主模板已定义）
%      \notecourseen 眉题英文（缺省 Course Notes）
%      \notesubtitle 副标题（缺省「学习笔记」）
%      \noteextra    底部附加信息（缺省不显示）
% ============================================================

% ---- 达·芬奇《蒙娜丽莎》取色（暖褐晕涂 + 蓝灰远景） ----
\providecolor{mlSky}{RGB}{142,164,166}      % 远景天空：蓝灰（上）
\providecolor{mlSky2}{RGB}{198,196,178}     % 远景亮部：暖灰
\providecolor{mlWater}{RGB}{120,146,152}    % 远水：青灰
\providecolor{mlGround}{RGB}{128,104,72}    % 中景地面：褐土
\providecolor{mlSkin}{RGB}{230,200,164}     % 肤色
\providecolor{mlHair}{RGB}{56,42,34}        % 头发：深褐
\providecolor{mlRobe}{RGB}{92,68,44}        % 衣裙：深褐
\providecolor{mlSleeve}{RGB}{152,114,68}    % 衣袖：赭黄
\providecolor{mlInk}{RGB}{46,34,26}         % 墨褐（标题/五官）

% ---- 缺省字段（主文件已定义的同名命令不会被覆盖） ----
\providecommand{\notecourseen}{Course Notes}
\providecommand{\notesubtitle}{学\ 习\ 笔\ 记}
\providecommand{\noteextra}{}

\begin{titlepage}
	\begin{tikzpicture}[remember picture, overlay]
		\begin{scope}[shift={(current page.south west)}]

		% ============================================================
		%  1. 远景（上冷下暖的晕涂天空）
		% ============================================================
		\shade[top color=mlSky, bottom color=mlSky2] (0,0) rectangle (21,29.7);

		% ============================================================
		%  2. 远山（层叠折线，越远越淡）
		% ============================================================
		\fill[mlSky2, opacity=0.75]
			(0,23.4) -- (2.6,25.2) -- (4.4,24.0) -- (6.8,25.8) -- (9.0,24.2)
			-- (12.0,26.0) -- (14.6,23.8) -- (17.2,25.4) -- (21,23.6)
			-- (21,20.4) -- (0,20.4) -- cycle;
		\fill[mlSky, opacity=0.45]
			(0,20.6) -- (3.0,22.2) -- (5.2,20.8) -- (8.0,22.6) -- (11.0,20.6)
			-- (14.2,22.4) -- (17.6,20.4) -- (21,21.8) -- (21,18.4) -- (0,18.4) -- cycle;

		% ============================================================
		%  3. 蜿蜒的河流与中景地面
		% ============================================================
		\draw[mlWater, line width=5.0pt, opacity=0.55, line cap=round]
			(1.6,17.6) .. controls (5.0,17.0) and (6.4,15.2) .. (9.6,14.8)
			.. controls (12.8,14.4) and (14.6,15.6) .. (18.0,14.6);
		\fill[mlGround, opacity=0.50]
			plot[domain=0:21, samples=50, smooth] (\x, {13.0 + 1.1*sin(deg(0.48*\x + 20))})
			-- (21,0) -- (0,0) -- cycle;

		% ============================================================
		%  4. 人物：衣裙与身形
		% ============================================================
		\fill[mlRobe]
			(4.2,0) .. controls (5.6,4.5) and (7.4,7.6) .. (8.9,10.2)
			.. controls (9.6,11.2) and (11.4,11.2) .. (12.1,10.2)
			.. controls (13.6,7.6) and (15.4,4.5) .. (16.8,0) -- cycle;
		\fill[mlSleeve, opacity=0.85]
			(4.6,0) .. controls (6.0,3.6) and (7.4,6.2) .. (8.6,8.4)
			.. controls (7.6,7.0) and (6.2,3.4) .. (5.4,0) -- cycle;
		\fill[mlSleeve, opacity=0.85]
			(16.4,0) .. controls (15.0,3.6) and (13.6,6.2) .. (12.4,8.4)
			.. controls (13.4,7.0) and (14.8,3.4) .. (15.6,0) -- cycle;

		% ============================================================
		%  5. 颈与头
		% ============================================================
		\fill[mlSkin] (9.55,9.2) -- (11.45,9.2) -- (11.45,11.8) -- (9.55,11.8) -- cycle;
		\fill[mlSkin] (10.5,13.5) ellipse[x radius=2.30, y radius=3.00];

		% ============================================================
		%  6. 垂发（自头顶包覆两侧，垂至肩部）
		% ============================================================
		\fill[mlHair]
			(10.5,16.5)
			.. controls (7.8,16.2) and (7.3,13.2) .. (7.6,10.4)
			.. controls (7.8,9.0) and (8.7,9.0) .. (8.8,10.6)
			.. controls (9.0,13.2) and (9.2,15.0) .. (10.5,15.1)
			.. controls (11.8,15.0) and (12.0,13.2) .. (12.2,10.6)
			.. controls (12.3,9.0) and (13.2,9.0) .. (13.4,10.4)
			.. controls (13.7,13.2) and (13.2,16.2) .. (10.5,16.5) -- cycle;

		% ============================================================
		%  7. 五官（含蓄的微笑）
		% ============================================================
		\draw[mlInk, line width=1.1pt, opacity=0.55]
			(9.30,13.95) .. controls (9.62,14.15) and (9.98,14.15) .. (10.28,13.95);
		\draw[mlInk, line width=1.1pt, opacity=0.55]
			(10.72,13.95) .. controls (11.02,14.15) and (11.38,14.15) .. (11.70,13.95);
		\fill[mlInk, opacity=0.85] (9.68,13.32) ellipse[x radius=0.34, y radius=0.19];
		\fill[mlInk, opacity=0.85] (11.32,13.32) ellipse[x radius=0.34, y radius=0.19];
		\draw[mlInk, line width=1.0pt, opacity=0.5]
			(10.50,13.05) .. controls (10.40,12.45) and (10.62,12.22) .. (10.74,12.40);
		\draw[mlInk, line width=1.4pt, opacity=0.7]
			(9.86,11.72) .. controls (10.28,11.48) and (10.72,11.48) .. (11.14,11.72);
		\draw[mlInk, line width=0.8pt, opacity=0.28]
			(9.20,12.30) .. controls (9.05,11.60) and (9.30,11.05) .. (9.70,11.10);
		\draw[mlInk, line width=0.8pt, opacity=0.28]
			(11.80,12.30) .. controls (11.95,11.60) and (11.70,11.05) .. (11.30,11.10);

		% ============================================================
		%  8. 交叠的双手
		% ============================================================
		\fill[mlSkin]
			(8.6,6.0) .. controls (9.4,5.1) and (11.0,4.9) .. (12.2,5.5)
			.. controls (12.9,5.9) and (12.7,6.7) .. (11.8,6.9)
			.. controls (10.6,7.1) and (9.2,7.0) .. (8.6,6.4) -- cycle;
		\draw[mlInk, opacity=0.35, line width=1.0pt] (9.70,5.56) -- (10.34,5.46);
		\draw[mlInk, opacity=0.35, line width=1.0pt] (10.86,5.62) -- (11.50,5.76);

		% ============================================================
		%  9. 标题衬底（一层薄雾般的暖白，保证文字可读）
		% ============================================================
		\fill[mlSky2, opacity=0.40]
			(2.6,16.9) -- (17.4,17.3) -- (18.0,22.6) -- (16.4,23.6) -- (2.2,22.9) -- cycle;

		% ============================================================
		% 10. 金色细边框
		% ============================================================
		\draw[gold, line width=0.7pt, opacity=0.6]
			([shift={(0.85,0.85)}]current page.south west) rectangle
			([shift={(-0.85,-0.85)}]current page.north east);

		% ============================================================
		% 11. 顶部眉题（课程英文小字 + 自适应金色短线）
		% ============================================================
		\node[anchor=north, text=mlInk, align=center]
			(mlEyebrow) at ([yshift=-3.4cm]current page.north) {%
			{\sffamily\fontsize{12}{15}\selectfont \uppercase{\notecourseen}}
		};
		\draw[gold, line width=1.0pt] ($(mlEyebrow.west)+(-0.25cm,0)$) -- ++(-1.1cm,0);
		\draw[gold, line width=1.0pt] ($(mlEyebrow.east)+(0.25cm,0)$) -- ++(1.1cm,0);

		% ============================================================
		% 12. 主标题（居中，使用 \notename）
		% ============================================================
		\node[anchor=center, align=center, text=mlInk]
			at ([yshift=-8.3cm]current page.north) {%
			{\fontsize{38}{48}\sffamily\bfseries \notename}
		};

		% ============================================================
		% 13. 金色菱形 + 双细线
		% ============================================================
		\coordinate (C) at ([yshift=-11.6cm]current page.north);
		\draw[gold, line width=1.1pt] (C) ++(-2.2cm,0) -- ++(1.55cm,0);
		\draw[gold, line width=1.1pt] (C) ++(0.65cm,0) -- ++(1.55cm,0);
		\node[fill=gold, minimum size=0.24cm, inner sep=0pt, rotate=45] at (C) {};

		% ============================================================
		% 14. 副标题（使用 \notesubtitle）
		% ============================================================
		\node[anchor=north, text=mlInk, align=center]
			at ([yshift=-12.4cm]current page.north) {
			{\fontsize{15}{20}\sffamily \notesubtitle}
		};

		% ============================================================
		% 15. 底部作者、日期、附加信息
		% ============================================================
		\coordinate (B) at ([yshift=2.7cm]current page.south);
		\draw[gold, line width=0.8pt, opacity=0.85] (B) ++(-1.1cm,0.95cm) -- ++(2.2cm,0);
		\node[anchor=south, text=mlInk, align=center] at (B) {
			\begin{tabular}{c}
				{\fontsize{19}{24}\sffamily\bfseries \noteauthor} \\[0.35cm]
				{\large \sffamily \today}
				\ifstrempty{\noteextra}{}{\\[0.3cm]{\normalsize\sffamily \noteextra}}
			\end{tabular}
		};

		% ============================================================
		% 16. 右下角落款（画作信息）
		% ============================================================
		\node[anchor=south east, text=mlInk, opacity=0.7, font=\footnotesize\itshape]
			at ([shift={(-1.35,1.25)}]current page.south east)
			{Leonardo da Vinci, Mona Lisa, c.1503--1519};

		\end{scope}
	\end{tikzpicture}
\end{titlepage}

%
%  配色：ml* 前缀（达·芬奇取色）。本文件用 \providecolor 兜底，
%        单独复制到任何笔记本也能编译；想整体换调改这里即可。
%
%  注意：整幅画面绘于 \begin{scope}[shift={(current page.south west)}] 内，
%        所有坐标以「页面左下角为原点、单位 cm」书写；脱离该 scope
%        直接写 (x,y)，图形会落到页面外而不可见。
%
%  可用字段（在主文件导言区用 \newcommand 预定义即可覆盖缺省值）：
%      \notename     主标题（必填，主模板已定义）
%      \noteauthor   作者（必填，主模板已定义）
%      \notecourseen 眉题英文（缺省 Course Notes）
%      \notesubtitle 副标题（缺省「学习笔记」）
%      \noteextra    底部附加信息（缺省不显示）
% ============================================================

% ---- 达·芬奇《蒙娜丽莎》取色（暖褐晕涂 + 蓝灰远景） ----
\providecolor{mlSky}{RGB}{142,164,166}      % 远景天空：蓝灰（上）
\providecolor{mlSky2}{RGB}{198,196,178}     % 远景亮部：暖灰
\providecolor{mlWater}{RGB}{120,146,152}    % 远水：青灰
\providecolor{mlGround}{RGB}{128,104,72}    % 中景地面：褐土
\providecolor{mlSkin}{RGB}{230,200,164}     % 肤色
\providecolor{mlHair}{RGB}{56,42,34}        % 头发：深褐
\providecolor{mlRobe}{RGB}{92,68,44}        % 衣裙：深褐
\providecolor{mlSleeve}{RGB}{152,114,68}    % 衣袖：赭黄
\providecolor{mlInk}{RGB}{46,34,26}         % 墨褐（标题/五官）

% ---- 缺省字段（主文件已定义的同名命令不会被覆盖） ----
\providecommand{\notecourseen}{Course Notes}
\providecommand{\notesubtitle}{学\ 习\ 笔\ 记}
\providecommand{\noteextra}{}

\begin{titlepage}
	\begin{tikzpicture}[remember picture, overlay]
		\begin{scope}[shift={(current page.south west)}]

		% ============================================================
		%  1. 远景（上冷下暖的晕涂天空）
		% ============================================================
		\shade[top color=mlSky, bottom color=mlSky2] (0,0) rectangle (21,29.7);

		% ============================================================
		%  2. 远山（层叠折线，越远越淡）
		% ============================================================
		\fill[mlSky2, opacity=0.75]
			(0,23.4) -- (2.6,25.2) -- (4.4,24.0) -- (6.8,25.8) -- (9.0,24.2)
			-- (12.0,26.0) -- (14.6,23.8) -- (17.2,25.4) -- (21,23.6)
			-- (21,20.4) -- (0,20.4) -- cycle;
		\fill[mlSky, opacity=0.45]
			(0,20.6) -- (3.0,22.2) -- (5.2,20.8) -- (8.0,22.6) -- (11.0,20.6)
			-- (14.2,22.4) -- (17.6,20.4) -- (21,21.8) -- (21,18.4) -- (0,18.4) -- cycle;

		% ============================================================
		%  3. 蜿蜒的河流与中景地面
		% ============================================================
		\draw[mlWater, line width=5.0pt, opacity=0.55, line cap=round]
			(1.6,17.6) .. controls (5.0,17.0) and (6.4,15.2) .. (9.6,14.8)
			.. controls (12.8,14.4) and (14.6,15.6) .. (18.0,14.6);
		\fill[mlGround, opacity=0.50]
			plot[domain=0:21, samples=50, smooth] (\x, {13.0 + 1.1*sin(deg(0.48*\x + 20))})
			-- (21,0) -- (0,0) -- cycle;

		% ============================================================
		%  4. 人物：衣裙与身形
		% ============================================================
		\fill[mlRobe]
			(4.2,0) .. controls (5.6,4.5) and (7.4,7.6) .. (8.9,10.2)
			.. controls (9.6,11.2) and (11.4,11.2) .. (12.1,10.2)
			.. controls (13.6,7.6) and (15.4,4.5) .. (16.8,0) -- cycle;
		\fill[mlSleeve, opacity=0.85]
			(4.6,0) .. controls (6.0,3.6) and (7.4,6.2) .. (8.6,8.4)
			.. controls (7.6,7.0) and (6.2,3.4) .. (5.4,0) -- cycle;
		\fill[mlSleeve, opacity=0.85]
			(16.4,0) .. controls (15.0,3.6) and (13.6,6.2) .. (12.4,8.4)
			.. controls (13.4,7.0) and (14.8,3.4) .. (15.6,0) -- cycle;

		% ============================================================
		%  5. 颈与头
		% ============================================================
		\fill[mlSkin] (9.55,9.2) -- (11.45,9.2) -- (11.45,11.8) -- (9.55,11.8) -- cycle;
		\fill[mlSkin] (10.5,13.5) ellipse[x radius=2.30, y radius=3.00];

		% ============================================================
		%  6. 垂发（自头顶包覆两侧，垂至肩部）
		% ============================================================
		\fill[mlHair]
			(10.5,16.5)
			.. controls (7.8,16.2) and (7.3,13.2) .. (7.6,10.4)
			.. controls (7.8,9.0) and (8.7,9.0) .. (8.8,10.6)
			.. controls (9.0,13.2) and (9.2,15.0) .. (10.5,15.1)
			.. controls (11.8,15.0) and (12.0,13.2) .. (12.2,10.6)
			.. controls (12.3,9.0) and (13.2,9.0) .. (13.4,10.4)
			.. controls (13.7,13.2) and (13.2,16.2) .. (10.5,16.5) -- cycle;

		% ============================================================
		%  7. 五官（含蓄的微笑）
		% ============================================================
		\draw[mlInk, line width=1.1pt, opacity=0.55]
			(9.30,13.95) .. controls (9.62,14.15) and (9.98,14.15) .. (10.28,13.95);
		\draw[mlInk, line width=1.1pt, opacity=0.55]
			(10.72,13.95) .. controls (11.02,14.15) and (11.38,14.15) .. (11.70,13.95);
		\fill[mlInk, opacity=0.85] (9.68,13.32) ellipse[x radius=0.34, y radius=0.19];
		\fill[mlInk, opacity=0.85] (11.32,13.32) ellipse[x radius=0.34, y radius=0.19];
		\draw[mlInk, line width=1.0pt, opacity=0.5]
			(10.50,13.05) .. controls (10.40,12.45) and (10.62,12.22) .. (10.74,12.40);
		\draw[mlInk, line width=1.4pt, opacity=0.7]
			(9.86,11.72) .. controls (10.28,11.48) and (10.72,11.48) .. (11.14,11.72);
		\draw[mlInk, line width=0.8pt, opacity=0.28]
			(9.20,12.30) .. controls (9.05,11.60) and (9.30,11.05) .. (9.70,11.10);
		\draw[mlInk, line width=0.8pt, opacity=0.28]
			(11.80,12.30) .. controls (11.95,11.60) and (11.70,11.05) .. (11.30,11.10);

		% ============================================================
		%  8. 交叠的双手
		% ============================================================
		\fill[mlSkin]
			(8.6,6.0) .. controls (9.4,5.1) and (11.0,4.9) .. (12.2,5.5)
			.. controls (12.9,5.9) and (12.7,6.7) .. (11.8,6.9)
			.. controls (10.6,7.1) and (9.2,7.0) .. (8.6,6.4) -- cycle;
		\draw[mlInk, opacity=0.35, line width=1.0pt] (9.70,5.56) -- (10.34,5.46);
		\draw[mlInk, opacity=0.35, line width=1.0pt] (10.86,5.62) -- (11.50,5.76);

		% ============================================================
		%  9. 标题衬底（一层薄雾般的暖白，保证文字可读）
		% ============================================================
		\fill[mlSky2, opacity=0.40]
			(2.6,16.9) -- (17.4,17.3) -- (18.0,22.6) -- (16.4,23.6) -- (2.2,22.9) -- cycle;

		% ============================================================
		% 10. 金色细边框
		% ============================================================
		\draw[gold, line width=0.7pt, opacity=0.6]
			([shift={(0.85,0.85)}]current page.south west) rectangle
			([shift={(-0.85,-0.85)}]current page.north east);

		% ============================================================
		% 11. 顶部眉题（课程英文小字 + 自适应金色短线）
		% ============================================================
		\node[anchor=north, text=mlInk, align=center]
			(mlEyebrow) at ([yshift=-3.4cm]current page.north) {%
			{\sffamily\fontsize{12}{15}\selectfont \uppercase{\notecourseen}}
		};
		\draw[gold, line width=1.0pt] ($(mlEyebrow.west)+(-0.25cm,0)$) -- ++(-1.1cm,0);
		\draw[gold, line width=1.0pt] ($(mlEyebrow.east)+(0.25cm,0)$) -- ++(1.1cm,0);

		% ============================================================
		% 12. 主标题（居中，使用 \notename）
		% ============================================================
		\node[anchor=center, align=center, text=mlInk]
			at ([yshift=-8.3cm]current page.north) {%
			{\fontsize{38}{48}\sffamily\bfseries \notename}
		};

		% ============================================================
		% 13. 金色菱形 + 双细线
		% ============================================================
		\coordinate (C) at ([yshift=-11.6cm]current page.north);
		\draw[gold, line width=1.1pt] (C) ++(-2.2cm,0) -- ++(1.55cm,0);
		\draw[gold, line width=1.1pt] (C) ++(0.65cm,0) -- ++(1.55cm,0);
		\node[fill=gold, minimum size=0.24cm, inner sep=0pt, rotate=45] at (C) {};

		% ============================================================
		% 14. 副标题（使用 \notesubtitle）
		% ============================================================
		\node[anchor=north, text=mlInk, align=center]
			at ([yshift=-12.4cm]current page.north) {
			{\fontsize{15}{20}\sffamily \notesubtitle}
		};

		% ============================================================
		% 15. 底部作者、日期、附加信息
		% ============================================================
		\coordinate (B) at ([yshift=2.7cm]current page.south);
		\draw[gold, line width=0.8pt, opacity=0.85] (B) ++(-1.1cm,0.95cm) -- ++(2.2cm,0);
		\node[anchor=south, text=mlInk, align=center] at (B) {
			\begin{tabular}{c}
				{\fontsize{19}{24}\sffamily\bfseries \noteauthor} \\[0.35cm]
				{\large \sffamily \today}
				\ifstrempty{\noteextra}{}{\\[0.3cm]{\normalsize\sffamily \noteextra}}
			\end{tabular}
		};

		% ============================================================
		% 16. 右下角落款（画作信息）
		% ============================================================
		\node[anchor=south east, text=mlInk, opacity=0.7, font=\footnotesize\itshape]
			at ([shift={(-1.35,1.25)}]current page.south east)
			{Leonardo da Vinci, Mona Lisa, c.1503--1519};

		\end{scope}
	\end{tikzpicture}
\end{titlepage}

%
%  配色：ml* 前缀（达·芬奇取色）。本文件用 \providecolor 兜底，
%        单独复制到任何笔记本也能编译；想整体换调改这里即可。
%
%  注意：整幅画面绘于 \begin{scope}[shift={(current page.south west)}] 内，
%        所有坐标以「页面左下角为原点、单位 cm」书写；脱离该 scope
%        直接写 (x,y)，图形会落到页面外而不可见。
%
%  可用字段（在主文件导言区用 \newcommand 预定义即可覆盖缺省值）：
%      \notename     主标题（必填，主模板已定义）
%      \noteauthor   作者（必填，主模板已定义）
%      \notecourseen 眉题英文（缺省 Course Notes）
%      \notesubtitle 副标题（缺省「学习笔记」）
%      \noteextra    底部附加信息（缺省不显示）
% ============================================================

% ---- 达·芬奇《蒙娜丽莎》取色（暖褐晕涂 + 蓝灰远景） ----
\providecolor{mlSky}{RGB}{142,164,166}      % 远景天空：蓝灰（上）
\providecolor{mlSky2}{RGB}{198,196,178}     % 远景亮部：暖灰
\providecolor{mlWater}{RGB}{120,146,152}    % 远水：青灰
\providecolor{mlGround}{RGB}{128,104,72}    % 中景地面：褐土
\providecolor{mlSkin}{RGB}{230,200,164}     % 肤色
\providecolor{mlHair}{RGB}{56,42,34}        % 头发：深褐
\providecolor{mlRobe}{RGB}{92,68,44}        % 衣裙：深褐
\providecolor{mlSleeve}{RGB}{152,114,68}    % 衣袖：赭黄
\providecolor{mlInk}{RGB}{46,34,26}         % 墨褐（标题/五官）

% ---- 缺省字段（主文件已定义的同名命令不会被覆盖） ----
\providecommand{\notecourseen}{Course Notes}
\providecommand{\notesubtitle}{学\ 习\ 笔\ 记}
\providecommand{\noteextra}{}

\begin{titlepage}
	\begin{tikzpicture}[remember picture, overlay]
		\begin{scope}[shift={(current page.south west)}]

		% ============================================================
		%  1. 远景（上冷下暖的晕涂天空）
		% ============================================================
		\shade[top color=mlSky, bottom color=mlSky2] (0,0) rectangle (21,29.7);

		% ============================================================
		%  2. 远山（层叠折线，越远越淡）
		% ============================================================
		\fill[mlSky2, opacity=0.75]
			(0,23.4) -- (2.6,25.2) -- (4.4,24.0) -- (6.8,25.8) -- (9.0,24.2)
			-- (12.0,26.0) -- (14.6,23.8) -- (17.2,25.4) -- (21,23.6)
			-- (21,20.4) -- (0,20.4) -- cycle;
		\fill[mlSky, opacity=0.45]
			(0,20.6) -- (3.0,22.2) -- (5.2,20.8) -- (8.0,22.6) -- (11.0,20.6)
			-- (14.2,22.4) -- (17.6,20.4) -- (21,21.8) -- (21,18.4) -- (0,18.4) -- cycle;

		% ============================================================
		%  3. 蜿蜒的河流与中景地面
		% ============================================================
		\draw[mlWater, line width=5.0pt, opacity=0.55, line cap=round]
			(1.6,17.6) .. controls (5.0,17.0) and (6.4,15.2) .. (9.6,14.8)
			.. controls (12.8,14.4) and (14.6,15.6) .. (18.0,14.6);
		\fill[mlGround, opacity=0.50]
			plot[domain=0:21, samples=50, smooth] (\x, {13.0 + 1.1*sin(deg(0.48*\x + 20))})
			-- (21,0) -- (0,0) -- cycle;

		% ============================================================
		%  4. 人物：衣裙与身形
		% ============================================================
		\fill[mlRobe]
			(4.2,0) .. controls (5.6,4.5) and (7.4,7.6) .. (8.9,10.2)
			.. controls (9.6,11.2) and (11.4,11.2) .. (12.1,10.2)
			.. controls (13.6,7.6) and (15.4,4.5) .. (16.8,0) -- cycle;
		\fill[mlSleeve, opacity=0.85]
			(4.6,0) .. controls (6.0,3.6) and (7.4,6.2) .. (8.6,8.4)
			.. controls (7.6,7.0) and (6.2,3.4) .. (5.4,0) -- cycle;
		\fill[mlSleeve, opacity=0.85]
			(16.4,0) .. controls (15.0,3.6) and (13.6,6.2) .. (12.4,8.4)
			.. controls (13.4,7.0) and (14.8,3.4) .. (15.6,0) -- cycle;

		% ============================================================
		%  5. 颈与头
		% ============================================================
		\fill[mlSkin] (9.55,9.2) -- (11.45,9.2) -- (11.45,11.8) -- (9.55,11.8) -- cycle;
		\fill[mlSkin] (10.5,13.5) ellipse[x radius=2.30, y radius=3.00];

		% ============================================================
		%  6. 垂发（自头顶包覆两侧，垂至肩部）
		% ============================================================
		\fill[mlHair]
			(10.5,16.5)
			.. controls (7.8,16.2) and (7.3,13.2) .. (7.6,10.4)
			.. controls (7.8,9.0) and (8.7,9.0) .. (8.8,10.6)
			.. controls (9.0,13.2) and (9.2,15.0) .. (10.5,15.1)
			.. controls (11.8,15.0) and (12.0,13.2) .. (12.2,10.6)
			.. controls (12.3,9.0) and (13.2,9.0) .. (13.4,10.4)
			.. controls (13.7,13.2) and (13.2,16.2) .. (10.5,16.5) -- cycle;

		% ============================================================
		%  7. 五官（含蓄的微笑）
		% ============================================================
		\draw[mlInk, line width=1.1pt, opacity=0.55]
			(9.30,13.95) .. controls (9.62,14.15) and (9.98,14.15) .. (10.28,13.95);
		\draw[mlInk, line width=1.1pt, opacity=0.55]
			(10.72,13.95) .. controls (11.02,14.15) and (11.38,14.15) .. (11.70,13.95);
		\fill[mlInk, opacity=0.85] (9.68,13.32) ellipse[x radius=0.34, y radius=0.19];
		\fill[mlInk, opacity=0.85] (11.32,13.32) ellipse[x radius=0.34, y radius=0.19];
		\draw[mlInk, line width=1.0pt, opacity=0.5]
			(10.50,13.05) .. controls (10.40,12.45) and (10.62,12.22) .. (10.74,12.40);
		\draw[mlInk, line width=1.4pt, opacity=0.7]
			(9.86,11.72) .. controls (10.28,11.48) and (10.72,11.48) .. (11.14,11.72);
		\draw[mlInk, line width=0.8pt, opacity=0.28]
			(9.20,12.30) .. controls (9.05,11.60) and (9.30,11.05) .. (9.70,11.10);
		\draw[mlInk, line width=0.8pt, opacity=0.28]
			(11.80,12.30) .. controls (11.95,11.60) and (11.70,11.05) .. (11.30,11.10);

		% ============================================================
		%  8. 交叠的双手
		% ============================================================
		\fill[mlSkin]
			(8.6,6.0) .. controls (9.4,5.1) and (11.0,4.9) .. (12.2,5.5)
			.. controls (12.9,5.9) and (12.7,6.7) .. (11.8,6.9)
			.. controls (10.6,7.1) and (9.2,7.0) .. (8.6,6.4) -- cycle;
		\draw[mlInk, opacity=0.35, line width=1.0pt] (9.70,5.56) -- (10.34,5.46);
		\draw[mlInk, opacity=0.35, line width=1.0pt] (10.86,5.62) -- (11.50,5.76);

		% ============================================================
		%  9. 标题衬底（一层薄雾般的暖白，保证文字可读）
		% ============================================================
		\fill[mlSky2, opacity=0.40]
			(2.6,16.9) -- (17.4,17.3) -- (18.0,22.6) -- (16.4,23.6) -- (2.2,22.9) -- cycle;

		% ============================================================
		% 10. 金色细边框
		% ============================================================
		\draw[gold, line width=0.7pt, opacity=0.6]
			([shift={(0.85,0.85)}]current page.south west) rectangle
			([shift={(-0.85,-0.85)}]current page.north east);

		% ============================================================
		% 11. 顶部眉题（课程英文小字 + 自适应金色短线）
		% ============================================================
		\node[anchor=north, text=mlInk, align=center]
			(mlEyebrow) at ([yshift=-3.4cm]current page.north) {%
			{\sffamily\fontsize{12}{15}\selectfont \uppercase{\notecourseen}}
		};
		\draw[gold, line width=1.0pt] ($(mlEyebrow.west)+(-0.25cm,0)$) -- ++(-1.1cm,0);
		\draw[gold, line width=1.0pt] ($(mlEyebrow.east)+(0.25cm,0)$) -- ++(1.1cm,0);

		% ============================================================
		% 12. 主标题（居中，使用 \notename）
		% ============================================================
		\node[anchor=center, align=center, text=mlInk]
			at ([yshift=-8.3cm]current page.north) {%
			{\fontsize{38}{48}\sffamily\bfseries \notename}
		};

		% ============================================================
		% 13. 金色菱形 + 双细线
		% ============================================================
		\coordinate (C) at ([yshift=-11.6cm]current page.north);
		\draw[gold, line width=1.1pt] (C) ++(-2.2cm,0) -- ++(1.55cm,0);
		\draw[gold, line width=1.1pt] (C) ++(0.65cm,0) -- ++(1.55cm,0);
		\node[fill=gold, minimum size=0.24cm, inner sep=0pt, rotate=45] at (C) {};

		% ============================================================
		% 14. 副标题（使用 \notesubtitle）
		% ============================================================
		\node[anchor=north, text=mlInk, align=center]
			at ([yshift=-12.4cm]current page.north) {
			{\fontsize{15}{20}\sffamily \notesubtitle}
		};

		% ============================================================
		% 15. 底部作者、日期、附加信息
		% ============================================================
		\coordinate (B) at ([yshift=2.7cm]current page.south);
		\draw[gold, line width=0.8pt, opacity=0.85] (B) ++(-1.1cm,0.95cm) -- ++(2.2cm,0);
		\node[anchor=south, text=mlInk, align=center] at (B) {
			\begin{tabular}{c}
				{\fontsize{19}{24}\sffamily\bfseries \noteauthor} \\[0.35cm]
				{\large \sffamily \today}
				\ifstrempty{\noteextra}{}{\\[0.3cm]{\normalsize\sffamily \noteextra}}
			\end{tabular}
		};

		% ============================================================
		% 16. 右下角落款（画作信息）
		% ============================================================
		\node[anchor=south east, text=mlInk, opacity=0.7, font=\footnotesize\itshape]
			at ([shift={(-1.35,1.25)}]current page.south east)
			{Leonardo da Vinci, Mona Lisa, c.1503--1519};

		\end{scope}
	\end{tikzpicture}
\end{titlepage}

%
%  配色：ml* 前缀（达·芬奇取色）。本文件用 \providecolor 兜底，
%        单独复制到任何笔记本也能编译；想整体换调改这里即可。
%
%  注意：整幅画面绘于 \begin{scope}[shift={(current page.south west)}] 内，
%        所有坐标以「页面左下角为原点、单位 cm」书写；脱离该 scope
%        直接写 (x,y)，图形会落到页面外而不可见。
%
%  可用字段（在主文件导言区用 \newcommand 预定义即可覆盖缺省值）：
%      \notename     主标题（必填，主模板已定义）
%      \noteauthor   作者（必填，主模板已定义）
%      \notecourseen 眉题英文（缺省 Course Notes）
%      \notesubtitle 副标题（缺省「学习笔记」）
%      \noteextra    底部附加信息（缺省不显示）
% ============================================================

% ---- 达·芬奇《蒙娜丽莎》取色（暖褐晕涂 + 蓝灰远景） ----
\providecolor{mlSky}{RGB}{142,164,166}      % 远景天空：蓝灰（上）
\providecolor{mlSky2}{RGB}{198,196,178}     % 远景亮部：暖灰
\providecolor{mlWater}{RGB}{120,146,152}    % 远水：青灰
\providecolor{mlGround}{RGB}{128,104,72}    % 中景地面：褐土
\providecolor{mlSkin}{RGB}{230,200,164}     % 肤色
\providecolor{mlHair}{RGB}{56,42,34}        % 头发：深褐
\providecolor{mlRobe}{RGB}{92,68,44}        % 衣裙：深褐
\providecolor{mlSleeve}{RGB}{152,114,68}    % 衣袖：赭黄
\providecolor{mlInk}{RGB}{46,34,26}         % 墨褐（标题/五官）

% ---- 缺省字段（主文件已定义的同名命令不会被覆盖） ----
\providecommand{\notecourseen}{Course Notes}
\providecommand{\notesubtitle}{学\ 习\ 笔\ 记}
\providecommand{\noteextra}{}

\begin{titlepage}
	\begin{tikzpicture}[remember picture, overlay]
		\begin{scope}[shift={(current page.south west)}]

		% ============================================================
		%  1. 远景（上冷下暖的晕涂天空）
		% ============================================================
		\shade[top color=mlSky, bottom color=mlSky2] (0,0) rectangle (21,29.7);

		% ============================================================
		%  2. 远山（层叠折线，越远越淡）
		% ============================================================
		\fill[mlSky2, opacity=0.75]
			(0,23.4) -- (2.6,25.2) -- (4.4,24.0) -- (6.8,25.8) -- (9.0,24.2)
			-- (12.0,26.0) -- (14.6,23.8) -- (17.2,25.4) -- (21,23.6)
			-- (21,20.4) -- (0,20.4) -- cycle;
		\fill[mlSky, opacity=0.45]
			(0,20.6) -- (3.0,22.2) -- (5.2,20.8) -- (8.0,22.6) -- (11.0,20.6)
			-- (14.2,22.4) -- (17.6,20.4) -- (21,21.8) -- (21,18.4) -- (0,18.4) -- cycle;

		% ============================================================
		%  3. 蜿蜒的河流与中景地面
		% ============================================================
		\draw[mlWater, line width=5.0pt, opacity=0.55, line cap=round]
			(1.6,17.6) .. controls (5.0,17.0) and (6.4,15.2) .. (9.6,14.8)
			.. controls (12.8,14.4) and (14.6,15.6) .. (18.0,14.6);
		\fill[mlGround, opacity=0.50]
			plot[domain=0:21, samples=50, smooth] (\x, {13.0 + 1.1*sin(deg(0.48*\x + 20))})
			-- (21,0) -- (0,0) -- cycle;

		% ============================================================
		%  4. 人物：衣裙与身形
		% ============================================================
		\fill[mlRobe]
			(4.2,0) .. controls (5.6,4.5) and (7.4,7.6) .. (8.9,10.2)
			.. controls (9.6,11.2) and (11.4,11.2) .. (12.1,10.2)
			.. controls (13.6,7.6) and (15.4,4.5) .. (16.8,0) -- cycle;
		\fill[mlSleeve, opacity=0.85]
			(4.6,0) .. controls (6.0,3.6) and (7.4,6.2) .. (8.6,8.4)
			.. controls (7.6,7.0) and (6.2,3.4) .. (5.4,0) -- cycle;
		\fill[mlSleeve, opacity=0.85]
			(16.4,0) .. controls (15.0,3.6) and (13.6,6.2) .. (12.4,8.4)
			.. controls (13.4,7.0) and (14.8,3.4) .. (15.6,0) -- cycle;

		% ============================================================
		%  5. 颈与头
		% ============================================================
		\fill[mlSkin] (9.55,9.2) -- (11.45,9.2) -- (11.45,11.8) -- (9.55,11.8) -- cycle;
		\fill[mlSkin] (10.5,13.5) ellipse[x radius=2.30, y radius=3.00];

		% ============================================================
		%  6. 垂发（自头顶包覆两侧，垂至肩部）
		% ============================================================
		\fill[mlHair]
			(10.5,16.5)
			.. controls (7.8,16.2) and (7.3,13.2) .. (7.6,10.4)
			.. controls (7.8,9.0) and (8.7,9.0) .. (8.8,10.6)
			.. controls (9.0,13.2) and (9.2,15.0) .. (10.5,15.1)
			.. controls (11.8,15.0) and (12.0,13.2) .. (12.2,10.6)
			.. controls (12.3,9.0) and (13.2,9.0) .. (13.4,10.4)
			.. controls (13.7,13.2) and (13.2,16.2) .. (10.5,16.5) -- cycle;

		% ============================================================
		%  7. 五官（含蓄的微笑）
		% ============================================================
		\draw[mlInk, line width=1.1pt, opacity=0.55]
			(9.30,13.95) .. controls (9.62,14.15) and (9.98,14.15) .. (10.28,13.95);
		\draw[mlInk, line width=1.1pt, opacity=0.55]
			(10.72,13.95) .. controls (11.02,14.15) and (11.38,14.15) .. (11.70,13.95);
		\fill[mlInk, opacity=0.85] (9.68,13.32) ellipse[x radius=0.34, y radius=0.19];
		\fill[mlInk, opacity=0.85] (11.32,13.32) ellipse[x radius=0.34, y radius=0.19];
		\draw[mlInk, line width=1.0pt, opacity=0.5]
			(10.50,13.05) .. controls (10.40,12.45) and (10.62,12.22) .. (10.74,12.40);
		\draw[mlInk, line width=1.4pt, opacity=0.7]
			(9.86,11.72) .. controls (10.28,11.48) and (10.72,11.48) .. (11.14,11.72);
		\draw[mlInk, line width=0.8pt, opacity=0.28]
			(9.20,12.30) .. controls (9.05,11.60) and (9.30,11.05) .. (9.70,11.10);
		\draw[mlInk, line width=0.8pt, opacity=0.28]
			(11.80,12.30) .. controls (11.95,11.60) and (11.70,11.05) .. (11.30,11.10);

		% ============================================================
		%  8. 交叠的双手
		% ============================================================
		\fill[mlSkin]
			(8.6,6.0) .. controls (9.4,5.1) and (11.0,4.9) .. (12.2,5.5)
			.. controls (12.9,5.9) and (12.7,6.7) .. (11.8,6.9)
			.. controls (10.6,7.1) and (9.2,7.0) .. (8.6,6.4) -- cycle;
		\draw[mlInk, opacity=0.35, line width=1.0pt] (9.70,5.56) -- (10.34,5.46);
		\draw[mlInk, opacity=0.35, line width=1.0pt] (10.86,5.62) -- (11.50,5.76);

		% ============================================================
		%  9. 标题衬底（一层薄雾般的暖白，保证文字可读）
		% ============================================================
		\fill[mlSky2, opacity=0.40]
			(2.6,16.9) -- (17.4,17.3) -- (18.0,22.6) -- (16.4,23.6) -- (2.2,22.9) -- cycle;

		% ============================================================
		% 10. 金色细边框
		% ============================================================
		\draw[gold, line width=0.7pt, opacity=0.6]
			([shift={(0.85,0.85)}]current page.south west) rectangle
			([shift={(-0.85,-0.85)}]current page.north east);

		% ============================================================
		% 11. 顶部眉题（课程英文小字 + 自适应金色短线）
		% ============================================================
		\node[anchor=north, text=mlInk, align=center]
			(mlEyebrow) at ([yshift=-3.4cm]current page.north) {%
			{\sffamily\fontsize{12}{15}\selectfont \uppercase{\notecourseen}}
		};
		\draw[gold, line width=1.0pt] ($(mlEyebrow.west)+(-0.25cm,0)$) -- ++(-1.1cm,0);
		\draw[gold, line width=1.0pt] ($(mlEyebrow.east)+(0.25cm,0)$) -- ++(1.1cm,0);

		% ============================================================
		% 12. 主标题（居中，使用 \notename）
		% ============================================================
		\node[anchor=center, align=center, text=mlInk]
			at ([yshift=-8.3cm]current page.north) {%
			{\fontsize{38}{48}\sffamily\bfseries \notename}
		};

		% ============================================================
		% 13. 金色菱形 + 双细线
		% ============================================================
		\coordinate (C) at ([yshift=-11.6cm]current page.north);
		\draw[gold, line width=1.1pt] (C) ++(-2.2cm,0) -- ++(1.55cm,0);
		\draw[gold, line width=1.1pt] (C) ++(0.65cm,0) -- ++(1.55cm,0);
		\node[fill=gold, minimum size=0.24cm, inner sep=0pt, rotate=45] at (C) {};

		% ============================================================
		% 14. 副标题（使用 \notesubtitle）
		% ============================================================
		\node[anchor=north, text=mlInk, align=center]
			at ([yshift=-12.4cm]current page.north) {
			{\fontsize{15}{20}\sffamily \notesubtitle}
		};

		% ============================================================
		% 15. 底部作者、日期、附加信息
		% ============================================================
		\coordinate (B) at ([yshift=2.7cm]current page.south);
		\draw[gold, line width=0.8pt, opacity=0.85] (B) ++(-1.1cm,0.95cm) -- ++(2.2cm,0);
		\node[anchor=south, text=mlInk, align=center] at (B) {
			\begin{tabular}{c}
				{\fontsize{19}{24}\sffamily\bfseries \noteauthor} \\[0.35cm]
				{\large \sffamily \today}
				\ifstrempty{\noteextra}{}{\\[0.3cm]{\normalsize\sffamily \noteextra}}
			\end{tabular}
		};

		% ============================================================
		% 16. 右下角落款（画作信息）
		% ============================================================
		\node[anchor=south east, text=mlInk, opacity=0.7, font=\footnotesize\itshape]
			at ([shift={(-1.35,1.25)}]current page.south east)
			{Leonardo da Vinci, Mona Lisa, c.1503--1519};

		\end{scope}
	\end{tikzpicture}
\end{titlepage}
